\begin{table}[t]
\caption{Positioning against adjacent reviews and surveys. Prior work supplies the clinical and methodological context; this review adds non-cumulative capability boundaries and an orthogonal evidence analysis.}
\label{tab:review-positioning}
\centering
\small
\begin{tabularx}{\linewidth}{@{}p{0.28\linewidth}p{0.26\linewidth}X@{}}
\toprule
\textbf{Literature} & \textbf{Primary organizing axis} & \textbf{What this review adds} \\
\midrule
Medical and biomedical world-model position/review articles \citep{Qazi2025BeyondGenerativeAI,Saeed2026MedicalWorldModel,Liu2026MedicalWorldModelsReview,Wang2026TowardsBiomedicalWorldModels} & Capability ladders; coupled state, dynamics, and intervention roadmaps; and multiscale biomedical roles from data engines to planning substrates & A non-cumulative boundary contract plus corpus-wide adjudication of implemented capability, clinical facets, uncertainty, and public evidence. \\
Digital-twin surveys and consensus work \citep{Iqbal2025consensusstatementon,Katsoulakis2024Digitaltwinshealth,Ghatti2023DigitalTwinsHealthcare} & Patient-specific virtual counterparts, mechanistic modeling, and clinical ambition & A learned-model capability scale that prevents digital-twin language from implying individualized rollout or validated action response before those are evidenced. \\
Structured-EHR foundation-model reviews \citep{Guo2026Systematicreviewfoundation} & EHR pretraining, trajectory modeling, and clinical event prediction & A state--action distinction that separates predicting care-process tokens from simulating patient response under settable interventions. \\
This review & Capability scale $\times$ SATO-V evidence contract, with application and architecture as secondary lenses & Paper-level claim calibration that separates implemented capability from action validity, causal support, clinical setting, and validation strength. \\
\bottomrule
\end{tabularx}
\end{table}
